\documentclass[12pt,a4paper]{article}

\usepackage{fontspec}
\usepackage{polyglossia}
\usepackage{amsmath}
\usepackage{amssymb}
\usepackage{geometry}
\usepackage{xcolor}
\usepackage{hyperref}
\usepackage{booktabs}
\usepackage{array}
\usepackage{fancyhdr}

\setmainlanguage{russian}
\setmainfont{Times New Roman}
\setsansfont{Arial}
\setmonofont{Menlo}

\geometry{
  top=22mm,
  right=22mm,
  bottom=24mm,
  left=24mm
}

\definecolor{accent}{HTML}{1F4D78}
\definecolor{muted}{HTML}{555555}
\definecolor{soft}{HTML}{F2F5F8}

\hypersetup{
  colorlinks=true,
  linkcolor=accent,
  urlcolor=accent,
  pdftitle={Расчёт вектора состояния в МПС RaspTank},
  pdfauthor={RaspTank Server}
}

\pagestyle{fancy}
\fancyhf{}
\lhead{\small Расчёт вектора состояния в МПС}
\rhead{\small RaspTank}
\cfoot{\small\thepage}
\renewcommand{\headrulewidth}{0.4pt}

\setlength{\parindent}{0pt}
\setlength{\parskip}{7pt}
\setlength{\headheight}{14pt}

\newcommand{\vect}[1]{\mathbf{#1}}

\begin{document}

\begin{center}
  {\LARGE\bfseries Расчёт вектора состояния в МПС}\\[4pt]
  {\large Модель пространства состояний и LQR-регулятор RaspTank}\\[10pt]
  {\color{muted}Технический фрагмент для отчёта}
\end{center}

\vspace{8pt}

\section{Вектор состояния}

В контуре МПС различаются три связанных вектора:

\[
\vect{x}_{\mathrm{real}}
=
\begin{bmatrix}
x_{\mathrm{real}} \\
y_{\mathrm{real}} \\
\theta_{\mathrm{real}} \\
v_{\mathrm{real}} \\
\omega_{\mathrm{real}}
\end{bmatrix},
\qquad
\vect{x}_{*}
=
\begin{bmatrix}
x_{*} \\
y_{*} \\
\theta_{*} \\
v_{*} \\
\omega_{*}
\end{bmatrix}.
\]

Вектор, который подаётся в LQR-регулятор, является ошибкой относительно
желаемого состояния:

\[
\vect{x}
=
\vect{x}_{*} - \vect{x}_{\mathrm{real}}
=
\begin{bmatrix}
x_{*} - x_{\mathrm{real}} \\
y_{*} - y_{\mathrm{real}} \\
\theta_{*} - \theta_{\mathrm{real}} \\
v_{*} - v_{\mathrm{real}} \\
\omega_{*} - \omega_{\mathrm{real}}
\end{bmatrix}.
\]

Обозначения:

\begin{itemize}
  \item \(\vect{x}_{\mathrm{real}}\) — реальное состояние машинки из оценки L2;
  \item \(\vect{x}_{*}\) — желаемое состояние, заданное командой движения;
  \item \(\vect{x}\) — ошибка состояния, используемая регулятором.
\end{itemize}

Реальное состояние берётся из датчиков и кинематики: \(x_{\mathrm{real}}\),
\(y_{\mathrm{real}}\), \(\theta_{\mathrm{real}}\) и \(v_{\mathrm{real}}\)
ведёт оцениватель положения L2, а \(\omega_{\mathrm{real}}\) берётся с
гироскопа при наличии измерения.

Желаемое состояние строится из активной целевой точки L3, но координата
\(x_{*}\) не берётся как фиксированное значение цели. Сначала целевая точка
обозначается как
\[
X_{\mathrm{goal}} =
\begin{bmatrix}
x_{\mathrm{goal}} \\
y_{\mathrm{goal}}
\end{bmatrix}.
\]

По текущему реальному положению вычисляются параметры прямой до цели:
\[
a = y_{\mathrm{goal}} - y_{\mathrm{real}},
\qquad
b = x_{\mathrm{goal}} - x_{\mathrm{real}}.
\]

Тогда желаемый курс вычисляется как
\[
\theta_{*} = \operatorname{atan2}(a, b).
\]

В коде используется именно \(\operatorname{atan2}\), а не простое
\(\arctan(a/b)\), потому что эта функция корректно обрабатывает \(b = 0\)
и сохраняет правильный квадрант целевой точки.

Для построения \(X_{*}\) используется формула прямой из скриншота:
\[
y_{*} = y_{\mathrm{goal}},
\qquad
k = \frac{a}{b},
\qquad
x_{*} = \frac{y_{*}}{k}
= \frac{y_{\mathrm{goal}}}{a/b}.
\]

То есть \(x_{*}\) рассчитывается из реальных текущих параметров \(a\) и
\(b\), а не подставляется напрямую как \(x_{\mathrm{goal}}\). Это делает
вектор ошибки
\[
\vect{x} = \vect{x}_{*} - \vect{x}_{\mathrm{real}}
\]
согласованным с прямой, по которой машинка должна выходить к цели.

Для вырожденных случаев, когда \(|a|\) или \(|b|\) практически равны нулю,
код не делит на ноль и использует исходные координаты цели:
\[
x_{*} = x_{\mathrm{goal}},
\qquad
y_{*} = y_{\mathrm{goal}}.
\]

Для остальных компонент вектора из скриншота задаются
\[
v_{*} = v_{\mathrm{des}},
\qquad
\omega_{*} = 0.
\]

Важно не смешивать два разных параметра скорости. \(v_{\mathrm{des}}\)
является желаемой линейной скоростью движения, а \(v_0\) является
номинальной скоростью линеаризации модели и входит в матрицу \(A\).
Если команда явно передаёт \(v_0\), матрица LQR считается по нему; иначе
используется \(v_{\mathrm{des}}\). Для защиты от вырожденной линеаризации
при малой скорости применяется нижняя граница \(|v_0| \ge 1\ \text{см/с}\).

\section{Непрерывная модель МПС}

Непрерывная модель пространства состояний записывается в виде

\[
\dot{\vect{x}}_{\mathrm{real}} = A\vect{x}_{\mathrm{real}} + B\vect{u}_{\mathrm{final}},
\]

где управляющий вектор равен

\[
\vect{u}_{\mathrm{final}}
=
\begin{bmatrix}
v_{\mathrm{cmd}} \\
\omega_{\mathrm{cmd}}
\end{bmatrix}.
\]

Матрица состояния:

\[
A =
\begin{bmatrix}
0 & 0 & 0 & 1 & 0 \\
0 & 0 & v_0 & 0 & 0 \\
0 & 0 & 0 & 0 & 1 \\
0 & 0 & 0 & -\frac{1}{T_v} & 0 \\
0 & 0 & 0 & 0 & -\frac{1}{T_\omega}
\end{bmatrix}.
\]

Матрица управления:

\[
B =
\begin{bmatrix}
0 & 0 \\
0 & 0 \\
0 & 0 \\
\frac{1}{T_v} & 0 \\
0 & \frac{1}{T_\omega}
\end{bmatrix}.
\]

Здесь \(T_v\) — постоянная времени линейной скорости,
\(T_\omega\) — постоянная времени угловой скорости,
\(v_0\) — номинальная скорость, относительно которой линеаризуется модель.

\section{Расшифровка выражения
\texorpdfstring{\(A\vect{x}_{\mathrm{real}} + B\vect{u}_{\mathrm{final}}\)}
{A x real + B u final}}

Подставим матрицы \(A\), \(B\), реальный вектор состояния
\(\vect{x}_{\mathrm{real}}\) и финальное управление
\(\vect{u}_{\mathrm{final}}\):

\[
\dot{\vect{x}}_{\mathrm{real}}
=
\begin{bmatrix}
\dot{x}_{\mathrm{real}} \\
\dot{y}_{\mathrm{real}} \\
\dot{\theta}_{\mathrm{real}} \\
\dot{v}_{\mathrm{real}} \\
\dot{\omega}_{\mathrm{real}}
\end{bmatrix}
=
\begin{bmatrix}
v_{\mathrm{real}} \\
v_0\theta_{\mathrm{real}} \\
\omega_{\mathrm{real}} \\
-\frac{1}{T_v}v_{\mathrm{real}} + \frac{1}{T_v}v_{\mathrm{cmd}} \\
-\frac{1}{T_\omega}\omega_{\mathrm{real}} + \frac{1}{T_\omega}\omega_{\mathrm{cmd}}
\end{bmatrix}.
\]

Отсюда получаются скалярные уравнения:

\[
\dot{x}_{\mathrm{real}} = v_{\mathrm{real}},
\]

\[
\dot{y}_{\mathrm{real}} = v_0\theta_{\mathrm{real}},
\]

\[
\dot{\theta}_{\mathrm{real}} = \omega_{\mathrm{real}},
\]

\[
\dot{v}_{\mathrm{real}} = \frac{v_{\mathrm{cmd}} - v_{\mathrm{real}}}{T_v},
\]

\[
\dot{\omega}_{\mathrm{real}} =
\frac{\omega_{\mathrm{cmd}} - \omega_{\mathrm{real}}}{T_\omega}.
\]

Первые три уравнения описывают кинематику движения, а последние два —
инерционность изменения линейной и угловой скоростей.

\section{Дискретное обновление состояния}

В программе состояние робота обновляется дискретно с шагом \(\Delta t\).
Сначала вычисляется новый курс:

\[
\theta_{k+1} = \theta_k + \omega_k \Delta t.
\]

Для более устойчивого обновления координат используется средний угол за шаг:

\[
\bar{\theta}_k =
\frac{\theta_k + \theta_{k+1}}{2}.
\]

После этого координаты обновляются по формулам:

\[
x_{k+1} =
x_k + v_k \cos(\bar{\theta}_k)\Delta t,
\]

\[
y_{k+1} =
y_k + v_k \sin(\bar{\theta}_k)\Delta t.
\]

Итоговый вектор состояния после шага:

\[
\vect{x}_{k+1}
=
\begin{bmatrix}
x_{k+1} \\
y_{k+1} \\
\theta_{k+1} \\
v_k \\
\omega_k
\end{bmatrix}.
\]

\section{Связь скоростей корпуса и гусениц}

Для дифференциальной платформы линейная и угловая скорости корпуса связаны
со скоростями правой и левой гусениц:

\[
v = \frac{v_R + v_L}{2},
\]

\[
\omega = \frac{v_R - v_L}{L},
\]

где \(L\) — расстояние между гусеницами.

Обратное преобразование, используемое при выдаче команды на моторы:

\[
v_R = v + \omega \frac{L}{2},
\]

\[
v_L = v - \omega \frac{L}{2}.
\]

\section{LQR-регулятор}

Для расчёта корректирующего управления используется закон

\[
\vect{u}_{\mathrm{corr}} = K\vect{x},
\]

где \(K\) — матрица коэффициентов LQR, а \(\vect{x}\) — вектор ошибки
\(\vect{x}_{*} - \vect{x}_{\mathrm{real}}\).

В текущей реализации вектор ошибки имеет вид

\[
\vect{x}
=
\begin{bmatrix}
x_{\mathrm{err}} \\
y_{\mathrm{err}} \\
\theta_{\mathrm{err}} \\
v_{\mathrm{err}} \\
\omega_{\mathrm{err}}
\end{bmatrix}.
\]

Позиционные ошибки берутся из активной целевой точки:

\[
x_{\mathrm{err}} = x_{*} - x_{\mathrm{real}},
\qquad
y_{\mathrm{err}} = y_{*} - y_{\mathrm{real}}.
\]

Ошибка курса:

\[
\theta_{\mathrm{err}} = \theta_{*} - \theta_{\mathrm{real}}.
\]

Ошибка линейной скорости:

\[
v_{\mathrm{err}} = v_{*} - v_{\mathrm{real}}.
\]

Ошибка угловой скорости:

\[
\omega_{\mathrm{err}} = \omega_{*} - \omega_{\mathrm{real}}.
\]

Поэтому фактически регулятор вычисляет

\[
\begin{bmatrix}
\Delta v \\
\Delta \omega
\end{bmatrix}
=
K
\begin{bmatrix}
x_{\mathrm{err}} \\
y_{\mathrm{err}} \\
\theta_{\mathrm{err}} \\
v_{\mathrm{err}} \\
\omega_{\mathrm{err}}
\end{bmatrix}.
\]

Коррекция линейной скорости ограничивается сверху по модулю настройкой
\(\Delta v_{\max}\), но затем финальная линейная команда дополнительно
зажимается желаемой скоростью. Это не даёт LQR разгонять машину быстрее,
чем задано пользователем:

\[
\Delta v_{\mathrm{lim}}
=
\operatorname{clip}(\Delta v,\ -\Delta v_{\max},\ \Delta v_{\max}).
\]

Для движения вперёд:

\[
v_{\mathrm{final}}
=
\operatorname{clip}(v_{\mathrm{des}} + \Delta v_{\mathrm{lim}},\ 0,\ v_{\mathrm{des}}),
\qquad v_{\mathrm{des}} > 0.
\]

Для движения назад:

\[
v_{\mathrm{final}}
=
\operatorname{clip}(v_{\mathrm{des}} + \Delta v_{\mathrm{lim}},\ v_{\mathrm{des}},\ 0),
\qquad v_{\mathrm{des}} < 0.
\]

Для нулевой желаемой скорости:

\[
v_{\mathrm{final}} = 0,
\qquad v_{\mathrm{des}} = 0.
\]

\[
\omega_{\mathrm{final}} = \omega_{\mathrm{des}} + \Delta \omega.
\]

\section{Вывод}

Верхняя формула дашборда \(\dot{\vect{x}} = A\vect{x} + B\vect{u}\)
показывает модель изменения физического состояния робота:
\[
\vect{x}_{\mathrm{real}} =
[x_{\mathrm{real}},\ y_{\mathrm{real}},\ \theta_{\mathrm{real}},\
v_{\mathrm{real}},\ \omega_{\mathrm{real}}]^T.
\]

В LQR-управлении используется вектор ошибки
\[
\vect{x} =
\vect{x}_{*} - \vect{x}_{\mathrm{real}} =
[x_{\mathrm{err}},\ y_{\mathrm{err}},\ \theta_{\mathrm{err}},\ v_{\mathrm{err}},\
\omega_{\mathrm{err}}]^T.
\]

В дашборде эта же связь выводится численно как
\[
\Delta X = X_{*} - X_{\mathrm{real}},
\]
то есть рядом показываются значения \(X_{*}\), \(X_{\mathrm{real}}\) и
получившаяся разность \(\Delta X\), а не только обозначения векторов.

Следовательно, текущая МПС использует целевую точку для ошибок \(x\), \(y\)
и \(\theta\), а датчики и кинематику — для реального состояния
\(\vect{x}_{\mathrm{real}}\).

\end{document}
